\chapter{问题定义与整体框架}
\label{chap:framework}

本章首先对多轮对话检索中的查询重写任务进行形式化定义，统一后续章节的符号与评估方式（第~\ref{sec:formal}~节），进而给出本文"历史剪枝—渐进式重写—检索反馈优化"三阶段查询重写框架（第~\ref{sec:framework}~节），并讨论训练与推理的一致性（第~\ref{sec:pipeline}~节）。

\section{任务形式化定义}
\label{sec:formal}

\subsection{符号约定}

给定大规模文档集合 $\mathcal{D}=\{d_1, d_2, \dots, d_{|\mathcal{D}|}\}$，一段多轮对话由若干轮次构成，记作
\begin{equation}
C = \{(q_1, a_1), (q_2, a_2), \dots, (q_n, a_n)\},
\end{equation}
其中 $q_t$ 为第 $t$ 轮用户查询，$a_t$ 为系统在第 $t$ 轮的回复文本。评估第 $i$ 轮查询时，记对话历史为
\begin{equation}
H_i = \{(q_1, a_1), (q_2, a_2), \dots, (q_{i-1}, a_{i-1})\}.
\end{equation}
将对话历史中被标注为"当前轮次金标证据"的文档集合记作 $G_i \subseteq \mathcal{D}$。

\subsection{多轮对话检索任务}

多轮对话检索任务要求学习一个检索函数 $\mathcal{R}$，使其在给定对话历史 $H_i$ 与当前查询 $q_i$ 条件下，从语料库 $\mathcal{D}$ 中召回一个按相关性排序的文档列表 $\mathcal{R}(H_i, q_i) = (\hat{d}_1, \hat{d}_2, \dots, \hat{d}_K)$，在 Recall@k、NDCG@k~\citep{jarvelin2002cumulated}、MRR@k 等信息检索指标上最大化 $\hat{d}$ 与金标 $G_i$ 之间的一致性。

\subsection{查询重写子任务}

本文聚焦于查询重写子任务：设计一个重写函数 $\mathcal{F}_\theta:(H_i, q_i) \mapsto q_i^\star$，使得改写后的独立查询 $q_i^\star$ 可直接输入任意现成的单轮检索器 $\mathcal{R}_0$（输入仅为单个查询）并取得最优检索性能。形式化地，该子任务的优化目标为：
\begin{equation}
\theta^\star = \arg\max_{\theta}\ \mathbb{E}_{(H_i, q_i)\sim \mathcal{C}}\Big[\text{Recall@}k\big(\mathcal{R}_0(\mathcal{F}_\theta(H_i, q_i)),\ G_i\big)\Big].
\label{eq:qr_obj}
\end{equation}

式~\eqref{eq:qr_obj} 清楚地表明：\textbf{重写器的训练目标应当直接对齐下游检索器的召回指标}，而非对齐人工书写的参考重写。这正是本文第~\ref{chap:rfr}~章 RFR 方案的核心设计原则。

\subsection{评估指标}
\label{sec:metric}

本文从三个维度评估重写质量：

\textbf{检索质量}：以下游检索器的 Recall@k（$k\in\{5,20\}$）、MRR@20、NDCG@20 作为主要指标，直接衡量重写对终端目标的影响。

\textbf{重写表层质量}：以 BLEU-4~\citep{papineni2002bleu}、ROUGE-1/L~\citep{lin2004rouge} 衡量重写与人工参考之间的词汇重叠度。注意，该指标仅作为辅助参考，不作为优化目标。

\textbf{推理效率}：统计每轮重写的平均生成长度、每千轮次 tokens/s 吞吐、显存占用，用于评估实用性。

\section{本文整体框架}
\label{sec:framework}

\subsection{设计动机}

针对第~\ref{chap:introduction}~章所归纳的三大瓶颈，本文提出了由三个相对独立但协同工作的模块构成的查询重写框架，统称为 \textsc{PruneRewriteAlign}（简记 PRA）：

\begin{itemize}
    \item \textbf{HP-Pruner}：话题漂移感知的对话历史剪枝器，解决瓶颈一"历史噪声建模不足"；
    \item \textbf{PRW}：两阶段渐进式重写器，解决瓶颈二"异质现象耦合训练"；
    \item \textbf{RFR}：检索反馈驱动的偏好对齐器，解决瓶颈三"重写目标与检索偏好失配"。
\end{itemize}

三模块的协同关系如图~\ref{fig:framework_overall}~所示：HP-Pruner 对对话历史进行剪枝，得到压缩后的上下文 $\tilde H_i$；PRW 接收 $(\tilde H_i, q_i)$ 并完成两阶段重写，得到候选 $q_i^{\star,0}$；RFR 以检索反馈作为偏好信号，通过 DPO 对 PRW 进行轻量化微调，最终得到用于部署的重写器 $\mathcal{F}_{\theta^\star}$。

\begin{figure}[t]
\centering
% 占位图：学弟本地替换为自制框架图
\fbox{\parbox{0.9\linewidth}{\centering HP-Pruner $\rightarrow$ PRW (Coreference $\rightarrow$ Ellipsis) $\rightarrow$ RFR (DPO)\\{\small 对话历史 $H_i$ 经剪枝得到 $\tilde H_i$，渐进式重写得到 $q_i^{\star}$，并以下游检索器反馈进行偏好对齐。}}}
\bicaption{本文 \textsc{PruneRewriteAlign} 整体框架}{Overall architecture of the proposed PruneRewriteAlign framework}
\label{fig:framework_overall}
\end{figure}

\subsection{三模块的功能界面}

为保证方法的工程可替换性与分析的可解耦性，本文对每个模块的输入输出界面进行严格定义：

\textbf{HP-Pruner}：输入 $(H_i, q_i)$，输出剪枝后的历史 $\tilde H_i = \{(q_{j}, a_{j})\mid s_{j,i} \geq \tau\}$，其中 $s_{j,i}$ 为第 $j$ 轮历史与当前查询 $q_i$ 的预测相关性得分，$\tau$ 为阈值。详细设计见第~\ref{chap:hpprune}~章。

\textbf{PRW}：输入 $(\tilde H_i, q_i)$，依次输出第一阶段重写 $q_i^{(1)}$（仅完成指代消解）与第二阶段重写 $q_i^{(2)}$（完成省略补全），最终重写 $q_i^{\star,0} := q_i^{(2)}$。详细设计见第~\ref{chap:prw}~章。

\textbf{RFR}：输入 $(\tilde H_i, q_i)$ 与候选 $q_i^{\star,0}$，通过 DPO 对 PRW 的策略 $\pi_\theta$ 进行微调，得到对齐后的重写器 $\pi_{\theta^\star}$，在推理阶段直接输出最终重写 $q_i^\star$。详细设计见第~\ref{chap:rfr}~章。

\subsection{为什么是三阶段？}

本文选择"剪枝—重写—对齐"这一三阶段结构而非端到端单模型，主要基于以下三点考量：

（1）\textbf{可解耦的错误分析}：在工业部署中，端到端模型一旦出错，工程师很难定位故障发生在"上下文理解"、"指代消解"还是"目标对齐"的哪一环。三阶段设计允许分别评估和诊断。

（2）\textbf{模块独立优化与替换}：HP-Pruner 是轻量级分类器、PRW 是生成式重写器、RFR 是偏好对齐后的生成器；三者训练数据、计算量、参数规模差异巨大，分阶段训练效率显著高于联合训练。

（3）\textbf{可控的计算成本}：剪枝模块显著降低 PRW/RFR 的平均输入长度，而 DPO 本身是离线优化方案，使得整体训练与推理成本保持在中等算力范围内可落地。

\section{训练-推理流程与一致性}
\label{sec:pipeline}

本文所提框架的训练-推理流程如图~\ref{fig:pipeline_flow}~所示。训练阶段分为三个相对独立的步骤：

\begin{figure}[t]
\centering
\fbox{\parbox{0.9\linewidth}{\centering
\textbf{训练阶段}\\
Step 1: 基于 TopiOCQA/QReCC 话题切换标注训练 HP-Pruner\\
Step 2: 合成两阶段重写数据，用 LoRA 在 Qwen2.5-1.5B/3B/7B 上 SFT 得到 PRW\\
Step 3: 用下游检索器的 Recall 差异构造偏好对，DPO 微调得到 RFR\\[1ex]
\textbf{推理阶段}\\
$H_i, q_i \to \text{HP-Pruner} \to \tilde H_i \to \text{PRW/RFR} \to q_i^\star \to \mathcal{R}_0 \to \text{Top-}k$
}}
\bicaption{训练与推理流程}{Training and inference pipeline}
\label{fig:pipeline_flow}
\end{figure}

\subsection{训练-推理一致性}

为避免训练与推理分布偏移，本文在三个环节刻意保持一致性：

\begin{itemize}
    \item HP-Pruner 在训练阶段与推理阶段使用完全相同的上下文格式化模板；
    \item PRW 第二阶段训练时使用的历史均为"第一阶段已替换指代的历史"，与推理阶段的两阶段串联保持一致；
    \item RFR 构造偏好对时，使用的重写候选均来自 PRW 的推理输出而非外部 GPT-4 改写，避免风格偏移。
\end{itemize}

\section{本章小结}

本章对多轮对话检索中的查询重写子任务进行了形式化定义，给出了以 Recall@k 为主要目标的优化目标函数，进而提出了本文的三阶段框架 \textsc{PruneRewriteAlign}，并就三模块的功能界面、训练-推理一致性进行了讨论。接下来的三章将分别详述 HP-Pruner（第~\ref{chap:hpprune}~章）、PRW（第~\ref{chap:prw}~章）与 RFR（第~\ref{chap:rfr}~章）的具体设计、训练细节与分析结果。
